\section{AIME I 1987}
        \begin{enumerate}
        	\item (\textit{AIME I 1987}) An ordered pair $(m,n)$ of non-negative integers is called "simple" if the addition $m+n$ in base $10$ requires no carrying. Find the number of simple ordered pairs of non-negative integers that sum to $1492$. 


 
 \textbf{Answer: }300



\textbf{Solution}

Since no carrying over is allowed, the range of possible values of any digit of $m$ is from $0$ to the respective \href{https://artofproblemsolving.com/wiki/index.php?title=Digit}{digit} in $1492$ (the values of $n$ are then fixed). Thus, the number of \href{https://artofproblemsolving.com/wiki/index.php?title=Ordered\_pair}{ordered pairs} will be $(1 + 1)(4 + 1)(9 + 1)(2 + 1) = 2\cdot 5\cdot 10\cdot 3 = \boxed{300}$.

 


 
If you do not understand the above solution, consider this. For every \href{https://artofproblemsolving.com/wiki/index.php?title=Positive\_integer}{positive integer} $m$, there is only one \href{https://artofproblemsolving.com/wiki/index.php?title=Whole\_number}{whole number} $n$ that you can add to it to obtain the required sum. Also, the total number of non-negative integers that are smaller than or equal to an integer $a$ is $(a + 1)$ because there are $(a - 1)$ positive integers that are less than it, in addition to $0$ and itself.

 


	\item (\textit{AIME I 1987}) What is the largest possible distance between two points, one on the sphere of radius 19 with center $(-2,-10,5),$ and the other on the sphere of radius 87 with center $(12,8,-16)$?


 
 \textbf{Answer: }137



\textbf{Solution 1}

The distance between the two centers of the spheres can be determined via the \href{https://artofproblemsolving.com/wiki/index.php?title=Distance\_formula}{distance formula} in three dimensions: $\sqrt{(12 - (-2))^2 + (8 - (-10))^2 + (-16 - 5)^2} = \sqrt{14^2 + 18^2 + 21^2} = 31$. The largest possible distance would be the sum of the two radii and the distance between the two centers, making it $19 + 87 + 31 = \boxed{137}$.

 



\textbf{Solution 2}

Since you have a lot of time on the AIME, you could spend 3 hours drawing a paper 3-dimensional graph of the two circles. Then, you would guess and check two random points on the circles until you get the two farthest points, and you find the distance will be $\boxed{137}$. (However, this could be very inaccurate and may take up a lot of time, but very much recommended :)

 


	\item (\textit{AIME I 1987}) By a proper divisor of a natural number we mean a positive integral divisor other than 1 and the number itself. A natural number greater than 1 will be called "nice" if it is equal to the product of its distinct proper divisors. What is the sum of the first ten nice numbers?


 
 \textbf{Answer: }182



\textbf{Solution 1}

Let $p(n)$ denote the product of the distinct proper divisors of $n$. A number $n$ is \textit{nice} in one of two instances: 

 
            \begin{enumerate}
        \item It has exactly two distinct \href{https://artofproblemsolving.com/wiki/index.php?title=Prime}{prime} divisors.
            \begin{enumerate}[label=(\roman*)]
                \item If we let $n = pq$, where $p,q$ are the prime factors, then its proper divisors are $p$ and $q$, and $p = p \cdot q = n$.
            \end{enumerate}
            
        \item It is the cube of a prime number.
            \begin{enumerate}[label=(\roman*)]
                \item If we let $n=p^3$ with $p$ prime, then its proper divisors are $p$ and $p^2$, and $p = p \cdot p^2 =n$.
            \end{enumerate}
            
            \end{enumerate}
            We now show that the above are the only two cases. Suppose that another nice number existed that does not fall into one of these two categories. Then we can either express it in the form $n = pqr$ (with $p,q$ prime and $r > 1$) or $n = p^e$ (with $e \neq 3$). In the former case, it suffices to note that $p(n) \ge (pr) \cdot (qr) = pqr^2 > pqr = n$. 

 In the latter case, then $p(n) = p \cdot p^2 \cdots p^{(e-1)} = p^{(e-1)e/2}$.

 For $p(n) = n$, we need $p^{(e-1)e/2} = p^e \Longrightarrow e^2 - e = 2e \Longrightarrow$ $e = 0$ or $e = 3$.  

 Since $e \neq 3$, in the case $e = 0 \Longrightarrow n = 1$ does not work. 

 Thus, listing out the first ten numbers to fit this form, $2 \cdot 3 = 6,\ 2^3 = 8,\ 2 \cdot 5 = 10,$ $\ 2 \cdot 7 = 14,\ 3 \cdot 5 = 15,\ 3 \cdot 7 = 21,$ $\ 2 \cdot 11 = 22,\ 2 \cdot 13 = 26,$ $\ 3^3 = 27,\ 3 \cdot 11 = 33$.\href{https://artofproblemsolving.com/wiki/index.php?title=Sum}{Summing} these yields $\boxed{182}$.

 



\textbf{Solution 2}

Alternatively, we could note that $n$ is only nice when it only has two proper divisors, which, when multiplied, clearly yield $n$. We know that when the \href{https://artofproblemsolving.com/wiki/index.php?title=Prime\_factorization}{prime factorization} of $n = a_1^{b_1} \cdot a_2^{b_2} \cdot a_3^{b_3} . . . \cdot a_m^{b_m}$, the number of factors $f(n)$ of $n$ is 
 \[f(n) = (b_1 + 1)(b_2 +1)(b_3 +1) . . . (b_m +1).\]

 Since $n$ is nice, it may only have $4$ factors ($1$, $n$, $p$, and $q$). This means that $f(n) = 4$. The number $4$ can only be factored into $(2)(2)$ or $(4)(1)$, which means that either $b_1 = 1$ and $b_2 = 1$, or $b_1 = 3$. Therefore the only two cases are $n = pq$, or $n = p^3$.And then continue.

 



\textbf{Solution 3(Quick Brute Force)}

The only numbers that are "nice" will contain four factors, including one and itself. The other two are a "factor pair", which, when multiplied, give the original number. Testing for "nice" numbers, we get:6, 8, 10, 14, 15, 21, 22, 26, 27, and 33.Adding them up gives $\boxed{182}$. ~Sigma\_Bob

 


	\item (\textit{AIME I 1987}) Find the area of the region enclosed by the graph of $|x-60|+|y|=\left|\frac{x}{4}\right|.$

 \textbf{Answer: }480



\textbf{Solution 1}


                    \begin{figure}[H]
                        \centering
                        \includegraphics[width=0.5\textwidth]{Solution/AIME_I_1987_4_Solution_Figure_1.png}
                    \end{figure}
                    

 Since $|y|$ is \href{https://artofproblemsolving.com/wiki/index.php?title=Nonnegative}{nonnegative}, $\left|\frac{x}{4}\right| \ge |x - 60|$. Solving this gives us two equations: $\frac{x}{4} \ge x - 60\ \mathrm{and} \ -\frac{x}{4} \le x - 60$. Thus, $48 \le x \le 80$. The \href{https://artofproblemsolving.com/wiki/index.php?title=Maximum}{maximum} and \href{https://artofproblemsolving.com/wiki/index.php?title=Minimum}{minimum} y value is when $|x - 60| = 0$, which is when $x = 60$ and $y = \pm 15$. Since the graph is \href{https://artofproblemsolving.com/wiki/index.php?title=Symmetry}{symmetric} about the y-axis, we just need \href{https://artofproblemsolving.com/wiki/index.php?title=Casework}{casework} upon $x$. $\frac{x}{4} > 0$, so we break up the condition $|x-60|$:

 
            \begin{itemize}
        \item $x - 60 > 0$. Then $y = -\frac{3}{4}x+60$.
        \item $x - 60 < 0$. Then $y = \frac{5}{4}x-60$.
            \end{itemize}
            The area of the region enclosed by the graph is that of the quadrilateral defined by the points $(48,0),\ (60,15),\ (80,0), \ (60,-15)$. Breaking it up into triangles and solving or using the \href{https://artofproblemsolving.com/wiki/index.php?title=Shoelace\_Theorem}{Shoelace Theorem}, we get $2 \cdot \frac{1}{2}(80 - 48)(15) = \boxed{480}$.

 



\textbf{Solution 2}

Since $|y|$ is the only present $y$ "term" in this equation, we know that the area must be symmetrical about the x-axis.

 We'll consider the area when $y>0$ and we only consider the portion enclosed with $y=0$. Then, we'll double that area since the graph is symmetrical. 

 Now, let us remove the absolute values: 

 When $x\ge 60$: $x-60+y=0.25x$. This rearranges to $y=-0.75x+60$.

 When $0\le x<60$: $60-x+y=0.25$. So $y=1.25x-60$.

 When $x<0$: $60-x+y=-0.25x$. So $y=0.75x-60$.

 By simple sketching, we see the shape that looks like the image in Solution 1 (graph it out and you'll see). We see that the partial area we seek in this part is the triangle with the vertices $(60,15)$, $(48,0)$, $(80,0)$. This triangle has an area of $(80-48)*15*0.5=240$. 

 Simply double the area and we get $\boxed{480}$ as our final answer.~hastapasta

 


	\item (\textit{AIME I 1987}) Find $3x^2 y^2$ if $x$ and $y$ are integers such that $y^2 + 3x^2 y^2 = 30x^2 + 517$.


 
 \textbf{Answer: }588



\textbf{Solution 1}

If we move the $x^2$ term to the left side, it is factorable with \href{https://artofproblemsolving.com/wiki/index.php?title=SFFT}{Simon's Favorite Factoring Trick}:

 
 \[(3x^2 + 1)(y^2 - 10) = 517 - 10\]
$507$ is equal to $3 \cdot 13^2$. Since $x$ and $y$ are integers, $3x^2 + 1$ cannot equal a multiple of three. $169$ doesn't work either, so $3x^2 + 1 = 13$, and $x^2 = 4$. This leaves $y^2 - 10 = 39$, so $y^2 = 49$. Thus, $3x^2 y^2 = 3 \times 4 \times 49 = \boxed{588}$.

 



\textbf{Solution 2}

First factor out the $y^2$ term on the left side. 

 $(y^2)(1+3x^2) = 30x^2 + 517$ 

 Then divide both sides by $(1+3x^2)$ so we isolate the $y^2$ term.

 $y^2 = (30x^2 + 517)/(1+3x^2)$

 Now we long divide the right side to get

 $y^2 = (30x^2 + 517)/(1+3x^2)$

 Ok now we have a Diophantine equation to proceed we long divide the right-side of the equation

 $y^2 = 10 + 507/(3x^2 + 1)$

 Due to the fact that $x$ and $y$ are integers $507/(3x^2 + 1)$ must be a integer too.

 We thus prime factorize and list the factors of $507$
$507 = 13^2 \times 3$

 The factors of $507$ are $1, 3, 13, 39, 169, 507$

 Now we set the denominator equal to $1, 3, 13, 39, 169, 507$

 We see that $13$ works as $x = 2$ and thus $y = 7$

 So $3x^2 y^2 = 3(xy)^2 = 3(14^2) = 3(200 - 4) = 600-12 = \boxed{588}$. ~blankbox

 


	\item (\textit{AIME I 1987}) Rectangle $ABCD$ is divided into four parts of equal area by five segments as shown in the figure, where $XY = YB + BC + CZ = ZW = WD + DA + AX$, and $PQ$ is parallel to $AB$.  Find the length of $AB$ (in cm) if $BC = 19$ cm and $PQ = 87$ cm.


 
                    \begin{figure}[H]
                        \centering
                        \includegraphics[width=0.5\textwidth]{Problem/AIME_I_1987_6_Problem_Figure_1.png}
                    \end{figure}
                    

 \textbf{Answer: }193



\textbf{Solution 1}

Since $XY = WZ$, $PQ = PQ$ and the \href{https://artofproblemsolving.com/wiki/index.php?title=Area}{areas} of the \href{https://artofproblemsolving.com/wiki/index.php?title=Trapezoid}{trapezoids} $PQZW$ and $PQYX$ are the same, then the heights of the trapezoids are the same.  Thus both trapezoids have area $\frac{1}{2} \cdot \frac{19}{2}(XY + PQ) = \frac{19}{4}(XY + 87)$.  This number is also equal to one quarter the area of the entire rectangle, which is $\frac{19\cdot AB}{4}$, so we have $AB = XY + 87$.

 In addition, we see that the \href{https://artofproblemsolving.com/wiki/index.php?title=Perimeter}{perimeter} of the rectangle is $2AB + 38 = XA + AD + DW + WZ + ZC + CB + BY + YX = 4XY$, so $AB + 19 = 2XY$.

 Solving these two equations gives $AB = \boxed{193}$.

 


 



\textbf{Solution 2}

Let $YB=a$, $CZ=b$, $AX=c$, and $WD=d$. First we drop a perpendicular from $Q$ to a point $R$ on $BC$ so $QR=h$. Since $XY = WZ$ and $PQ = PQ$ and the \href{https://artofproblemsolving.com/wiki/index.php?title=Area}{areas} of the \href{https://artofproblemsolving.com/wiki/index.php?title=Trapezoid}{trapezoids} $PQZW$ and $PQYX$ are the same, the heights of the trapezoids are both $\frac{19}{2}$.From here, we have that $[BYQZC]=\frac{a+h}{2}*19/2+\frac{b+h}{2}*19/2=19/2* \frac{a+b+2h}{2}$. We are told that this area is equal to $[PXYQ]=\frac{19}{2}* \frac{XY+87}{2}=\frac{19}{2}* \frac{a+b+106}{2}$. Setting these equal to each other and solving gives $h=53$. In the same way, we find that the perpendicular from $P$ to $AD$ is $53$. So $AB=53*2+87=\boxed{193}$




\textbf{Solution 3}

Since $XY = YB + BC + CZ = ZW = WD + DA + AX$. Let $a = AX + DW = BY + CZ$. Since $2AB - 2a = XY = WZ$, then $XY = AB - a$.Let $S$ be the midpoint of $DA$, and $T$ be the midpoint of $CB$. Since the area of $PQWZ$ and $PQYX$ are the same, then their heights are the same, and so $PQ$ is \href{https://artofproblemsolving.com/wiki/index.php?title=Equidistant}{equidistant} from $AB$ and $CD$. This means that $PS$ is perpendicular to $DA$, and $QT$ is perpendicular to $BC$. Therefore, $PSCW$, $PSAX$, $QZCT$, and $QYTB$ are all trapezoids, and $QT = (AB - 87)/$2. This implies that 
 \[((a + 2((AB - 87)/2)/2) \cdot 19 = (((AB - a) + 87)/2) \cdot 19\] 
 \[(a + AB - 87) = AB - a + 87\] 
 \[2a = 174\] 
 \[a = 87\] Since $a + CB = XY$, $XY = 19 + 87 = 106$, and $AB = 106 + 87 = \boxed{193}$.

 


	\item (\textit{AIME I 1987}) Let $[r,s]$ denote the least common multiple of positive integers $r$ and $s$.  Find the number of ordered triples $(a,b,c)$ of positive integers for which $[a,b] = 1000$, $[b,c] = 2000$, and $[c,a] = 2000$.


 
 \textbf{Answer: }070



\textbf{Solution 1}

It's clear that we must have $a = 2^j5^k$, $b = 2^m 5^n$ and $c = 2^p5^q$ for some \href{https://artofproblemsolving.com/wiki/index.php?title=Nonnegative}{nonnegative} \href{https://artofproblemsolving.com/wiki/index.php?title=Integer}{integers} $j, k, m, n, p, q$.  Dealing first with the powers of 2: from the given conditions, $\max(j, m) = 3$, $\max(m, p) = \max(p, j) = 4$.  Thus we must have $p = 4$ and at least one of $m, j$ equal to 3.  This gives 7 possible triples $(j, m, p)$: $(0, 3, 4), (1, 3, 4), (2, 3, 4), (3, 3, 4), (3, 2, 4), (3, 1, 4)$ and $(3, 0, 4)$.

 Now, for the powers of 5: we have $\max(k, n) = \max(n, q) = \max(q, k) = 3$.  Thus, at least two of $k, n, q$ must be equal to 3, and the other can take any value between 0 and 3.  This gives us a total of 10 possible triples: $(3, 3, 3)$ and three possibilities of each of the forms $(3, 3, n)$, $(3, n, 3)$ and $(n, 3, 3)$.

 Since the \href{https://artofproblemsolving.com/wiki/index.php?title=Exponent}{exponents} of 2 and 5 must satisfy these conditions independently, we have a total of $7 \cdot 10 = 70$ possible valid triples.

 



\textbf{Solution 2}

$1000 = 2^35^3$ and $2000 = 2^45^3$. By \href{https://artofproblemsolving.com/wiki/index.php?title=LCM\#Using\_prime\_factorization}{looking at the prime factorization} of $2000$, $c$ must have a \href{https://artofproblemsolving.com/wiki/index.php?title=Factor}{factor} of $2^4$. If $c$ has a factor of $5^3$, then there are two \href{https://artofproblemsolving.com/wiki/index.php?title=Casework}{cases}: either (1) $a$ or $b = 5^32^3$, or (2) one of $a$ and $b$ has a factor of $5^3$ and the other a factor of $2^3$. For case 1, the other number will be in the form of $2^x5^y$, so there are $4 \cdot 4 = 16$ possible such numbers; since this can be either $a$ or $b$ there are a total of $2(16)-1=31$ possibilities. For case 2, $a$ and $b$ are in the form of $2^35^x$ and $2^y5^3$, with $x < 3$ and $y < 3$ (if they were equal to 3, it would overlap with case 1). Thus, there are $2(3 \cdot 3) = 18$ cases.

 If $c$ does not have a factor of $5^3$, then at least one of $a$ and $b$ must be $2^35^3$, and both must have a factor of $5^3$. Then, there are $4$ solutions possible just considering $a = 2^35^3$, and a total of $4 \cdot 2 - 1 = 7$ possibilities. Multiplying by three, as $0 \le c \le 2$, there are $7 \cdot 3 = 21$. Together, that makes $31 + 18 + 21 = 70$ solutions for $(a, b, c)$.

 


	\item (\textit{AIME I 1987}) What is the largest positive integer $n$ for which there is a unique integer $k$ such that $\frac{8}{15} < \frac{n}{n + k} < \frac{7}{13}$?


 
 \textbf{Answer: }112



\textbf{Solution 1}

Multiplying out all of the \href{https://artofproblemsolving.com/wiki/index.php?title=Denominator}{denominators}, we get:

 
 \begin{align*}104(n+k) &< 195n< 105(n+k)\\ 0 &< 91n - 104k < n + k\end{align*}


 Since $91n - 104k < n + k$, $k > \frac{6}{7}n$. Also, $0 < 91n - 104k$, so $k < \frac{7n}{8}$. Thus, $48n < 56k < 49n$. $k$ is unique if it is within a maximum \href{https://artofproblemsolving.com/wiki/index.php?title=Range}{range} of $112$, so $n = \boxed{112}$.

 



\textbf{Solution 2}

Flip all of the fractions for

 
 \[\begin{array}{ccccc}\frac{15}{8} &>& \frac{k + n}{n} &>& \frac{13}{7}\\  105n &>& 56 (k + n)& >& 104n\\  49n &>& 56k& >& 48n\end{array}\]

 Continue as in Solution 1.

 



\textbf{Solution 3}

Flip the fractions and subtract one from all sides to yield 
 \[\frac{7}{8}>\frac{k}{n}>\frac{6}{7}.\]  Multiply both sides by $56n$ to get 
 \[49n>56k>48n.\] This is equivalent to find the largest value of $n$ such that there is only one multiple of 56 within the open interval between $48n$ and $49n$.  If $n=112,$ then $98>k>96$ and $k=97$ is the unique value. For $n\geq 113,$ there is at least $(49\cdot 113-48\cdot 113)-1=112$ possible numbers for $k$ and there is one $k$ every 56 numbers. Hence, there must be at least two values of $k$ that work.  So, the largest value of $n$ is $\boxed{112}$.

 



\textbf{Solution 4}

Notice that in order for $k$ to be unique, $\frac{n}{n + k+ 1} \le \frac{8}{15}$ and $\frac{n}{n+ k-1} \ge \frac{7}{13}$ must be true. Solving these inequalities for $k$ yields $\frac{7}{6}(k-1) \le n \le \frac{8}{7}(k+1)$.

 Thus, we want to find $k$ such that $\frac{7}{6}(k-1)\le \frac{8}{7}(k+1)$. Solving this inequality yields $k \le 97$, and plugging this into $\frac{n}{n+k} < \frac{7}{13}$ in the original equation yields $n \le 112$ so the answer is $\boxed{112}$.

 


	\item (\textit{AIME I 1987}) Triangle $ABC$ has right angle at $B$, and contains a point $P$ for which $PA = 10$, $PB = 6$, and $\angle APB = \angle BPC = \angle CPA$.  Find $PC$.


 
                    \begin{figure}[H]
                        \centering
                        \includegraphics[width=0.5\textwidth]{Problem/AIME_I_1987_9_Problem_Figure_1.png}
                    \end{figure}
                    

 \textbf{Answer: }033



\textbf{Solution}

Let $PC = x$. Since $\angle APB = \angle BPC = \angle CPA$, each of them is equal to $120^\circ$.  By the \href{https://artofproblemsolving.com/wiki/index.php?title=Law\_of\_Cosines}{Law of Cosines} applied to triangles $\triangle APB$, $\triangle BPC$ and $\triangle CPA$ at their respective angles $P$, remembering that $\cos 120^\circ = -\frac12$, we have

 
 \[AB^2 = 36 + 100 + 60 = 196, BC^2 = 36 + x^2 + 6x, CA^2 = 100 + x^2 + 10x\]

 Then by the \href{https://artofproblemsolving.com/wiki/index.php?title=Pythagorean\_Theorem}{Pythagorean Theorem}, $AB^2 + BC^2 = CA^2$, so 

 
 \[x^2 + 10x + 100 = x^2 + 6x + 36 + 196\] 

 and

 
 \[4x = 132 \Longrightarrow x = \boxed{033}.\]

 \textbf{Note} This is the \href{https://artofproblemsolving.com/wiki/index.php?title=Fermat\_point}{Fermat point} of the triangle.

 


	\item (\textit{AIME I 1987}) Al walks down to the bottom of an escalator that is moving up and he counts 150 steps. His friend, Bob, walks up to the top of the escalator and counts 75 steps. If Al's speed of walking (in steps per unit time) is three times Bob's walking speed, how many steps are visible on the escalator at a given time? (Assume that this value is constant.) 


 
 \textbf{Answer: }120



\textbf{Solution 1}

Let the total number of steps be $x$, the speed of the escalator be $e$ and the speed of Bob be $b$.

 In the time it took Bob to climb up the escalator he saw 75 steps and also climbed the entire escalator.  Thus the contribution of the escalator must have been an additional $x - 75$ steps.  Since Bob and the escalator were both moving at a constant speed over the time it took Bob to climb, the \href{https://artofproblemsolving.com/wiki/index.php?title=Ratio}{ratio} of their distances covered is the same as the ratio of their speeds, so $\frac{e}{b} = \frac{x - 75}{75}$.  

 Similarly, in the time it took Al to walk down the escalator he saw 150 steps, so the escalator must have moved $150 - x$ steps in that time.  Thus $\frac{e}{3b} = \frac{150 - x}{150}$ or $\frac{e}{b} = \frac{150 - x}{50}$.

 Equating the two values of $\frac{e}{b}$ we have $\frac{x - 75}{75} = \frac{150 - x}{50}$ and so $2x - 150 = 450 - 3x$ and $5x = 600$ and $x = \boxed{120}$, the answer.

 



\textbf{Solution 2}

Again, let the total number of steps be $x$, the speed of the escalator be $e$ and the speed of Bob be $b$ (all "per unit time").

 Then this can be interpreted as a classic chasing problem: Bob is "behind" by $x$ steps, and since he moves at a pace of $b+e$ relative to the escalator, it will take $\frac{x}{b+e}=\frac{75}{e}$ time to get to the top.

 Similarly, Al will take $\frac{x}{3b-e}=\frac{150}{e}$ time to get to the bottom.

 From these two equations, we arrive at $150=\frac{ex}{3b-e}=2\cdot75=\frac{2ex}{b+e}=\frac{6ex}{3b+3e}=\frac{(ex)-(6ex)}{(3b-e)-(3b+3e)}=\frac{5x}{4}$$\implies600=5x\implies x=\boxed{120}$, where we have used the fact that $\frac{a}{b}=\frac{c}{d}=\frac{a\pm c}{b\pm d}$ (the proportion manipulations are motivated by the desire to isolate $x$, prompting the isolation of the $150$ on one side, and the fact that if we could cancel out the $b$'s, then the $e$'s in the numerator and denominator would cancel out, resulting in an equation with $x$ by itself).

 



\textbf{Solution 3}

Let $e$ and $b$ be the speeds of the escalator and Bob, respectively.

 When Al was on his way down, he took $150$ steps with a speed of $3b-e$ per step. When Bob was on his way up, he took $75$ steps with a speed of $b+e$ per step. Since Al and Bob were walking the same distance, we have
 \[\frac{150}{3b}(3b-e)=\frac{75}{b}(b+e)\]Solving gets the ratio $\frac{e}{b}=\frac{3}{5}$.

 Thus while Bob took $75$ steps to go up, the escalator contributed an extra $\frac{3}{5}\cdot75=45$ steps.

 Finally, there is a total of $75+45=\boxed{120}$ steps in the length of the escalator.

 



\textbf{Solution 4}

Please understand the machinery of an escalator before proceeding to read this solution.

 Let the number of steps that disappear at the top of the escalator equal $x.$ Assume that Al takes $3$ steps per second and that Bob takes $1$ step per second. Since Al counts $150$ steps, it takes him 
 \[\frac{150}{3}=50\] seconds to traverse the distance of the escalator moving downwards. Since Bob counts $75$ steps, it takes him 
 \[\frac{75}{1}=75\] seconds to traverse the distance of the escalator moving downwards.

 For the sake of this solution, we activate the emergency stop button on the escalator.

 Now, the escalator is not moving, or is simply a staircase. Imagine that Al is taking $3$ steps downwards every second, but we throw hands at him immediately after each second, such that he flinches and moves himself backwards $x$ steps. This is equivalent to Al taking $3-x$ steps downwards every second. Since we discovered that it takes him $50$ seconds to get from the top to the bottom of the escalator, and we are forcing Al to imitate the movement of the escalator, it also takes him $50$ seconds to move from the top to the bottom of the staircase. Thus, Al takes a total of 
 \[(3-x) \cdot 50=150-50x \qquad (\heartsuit)\] steps.

 The explanation for Bob is similar except now we pick him up and place him forward $x$ steps immediately after he takes his usual step per second, and since we discovered he does this for $75$ seconds, it takes him 
 \[(1+x) \cdot 75=75+75x \qquad (\clubsuit)\] steps to get from the bottom to the top. 

 
Note that because the escalator is broken and is now a staircase, Al and Bob must have had to take an equal amount of steps to get from the bottom to the top or from the top to the bottom. (Clearly, there are an equal amount of steps from the bottom to the top, and from the top to the bottom.) Therefore, we may equate $\heartsuit$ and $\clubsuit$ to get 
 \[150-50x=75+75x\] 
 \[x=\frac{3}{5}.\] Therefore, substituting $x$ in the expression we discovered in $\heartsuit,$ Al takes a total of 
 \[150-50x=150-50\left(\frac{3}{5}\right)=\boxed{120}\] steps, and we are done.

 



\textbf{Solution 5 (simple)}

WLOG, let Al's speed be $15$ steps per second, so Bob's speed is $5$ steps per second. Then, Al was on the escalator for $\frac{150}{15}\ = 10$ seconds and Bob was on the escalator for $\frac{75}{5}\ = 15$ seconds. Let $r$ be the rate of the escalator, in steps per second. Then, the total amount of steps is $150 - 10r = 75 + 15r$. Al is getting $10$ seconds of resistance at rate $r$ from the escalator, while Bob is getting $15$ seconds of help at rate $r$. Solving for $r$, we have $r = 3$ steps per second. Then, we can plug $r$ into the previous equation or subtract/add it to Al/Bob's rate (respectively) then multiply by their respective time. Either way, we get $\boxed{120}$ and we are done.

 


	\item (\textit{AIME I 1987}) Find the largest possible value of $k$ for which $3^{11}$ is expressible as the sum of $k$ consecutive positive integers.


 
 \textbf{Answer: }486



\textbf{Solution 1}

Let us write down one such sum, with $m$ terms and first term $n + 1$:

 $3^{11} = (n + 1) + (n + 2) + \ldots + (n + m) = \frac{1}{2} m(2n + m + 1)$.

 Thus $m(2n + m + 1) = 2 \cdot 3^{11}$ so $m$ is a \href{https://artofproblemsolving.com/wiki/index.php?title=Divisor}{divisor} of $2\cdot 3^{11}$.  However, because $n \geq 0$ we have $m^2 < m(m + 1) \leq 2\cdot 3^{11}$ so $m < \sqrt{2\cdot 3^{11}} < 3^6$.  Thus, we are looking for large factors of $2\cdot 3^{11}$ which are less than $3^6$.  The largest such factor is clearly $2\cdot 3^5 = 486$; for this value of $m$ we do indeed have the valid \href{https://artofproblemsolving.com/wiki/index.php?title=Expression}{expression} $3^{11} = 122 + 123 + \ldots + 607$, for which $k=\boxed{486}$.

 



\textbf{Solution 2}

First note that if $k$ is odd, and $n$ is the middle term, the sum equals $kn$. If $k$ is even, then we have the sum equal to $kn+k/2$, which will be even. Since $3^{11}$ is odd, we see that $k$ is odd. 

 Thus, we have $nk=3^{11} \implies n=3^{11}/k$. Also, note $n-(k+1)/2=0 \implies n=(k+1)/2.$ Subsituting $n=3^{11}/k$, we have $k^2+k=2*3^{11}$. Proceed as in solution 1.

 



\textbf{Solution 3}

Proceed as in Solution 1 until it is noted that $m$ is a divisor of $2\cdot 3^{11}$. The divisors of $2\cdot 3^{11}$ are $3^{1} , 2\cdot 3^{1} , 3^{2} , 2\cdot 3^{2} , \ldots , 2\cdot 3^{10} , 3^{11}$. Note that the factors of $m(2n + m + 1)$ are of opposite parity (if $m$ is odd, then $(2n + m + 1)$ is even and vice versa). Thus, one of the two factors will be a power of three, and the other will be twice a power of three. $(2n + m + 1)$ will represent the greater factor while $m$ will represent the lesser factor. Given this information, we need to find the factor pair that maximizes the lesser of the two factors, as this will maximize the value of $m$. The factor pair which maximizes the lesser factor is $2\cdot 3^{5}$ and $3^{6}$. It follows that $m$ = $2\cdot 3^{5}$ = $\boxed{486}$.

 


	\item (\textit{AIME I 1987}) Let $m$ be the smallest \href{https://artofproblemsolving.com/wiki/index.php?title=Integer}{integer} whose cube root is of the form $n+r$, where $n$ is a \href{https://artofproblemsolving.com/wiki/index.php?title=Positive\_integer}{positive integer} and $r$ is a positive real number less than $1/1000$. Find $n$.


 
 \textbf{Answer: }019



\textbf{Solution 1}

In order to keep $m$ as small as possible, we need to make $n$ as small as possible.

 $m = (n + r)^3 = n^3 + 3n^2r + 3nr^2 + r^3$.  Since $r < \frac{1}{1000}$ and $m - n^3 = r(3n^2 + 3nr + r^2)$ is an integer, we must have that $3n^2 + 3nr + r^2 \geq \frac{1}{r} > 1000$.  This means that the smallest possible $n$ should be less than 1000. In particular, $3nr + r^2$ should be less than 1, so $3n^2 > 999$ and $n > \sqrt{333}$.  $18^2 = 324 < 333 < 361 = 19^2$, so we must have $n \geq 19$.  Since we want to minimize $n$, we take $n = 19$.  Then for any positive value of $r$, $3n^2 + 3nr + r^2 > 3\cdot 19^2 > 1000$, so it is possible for $r$ to be less than $\frac{1}{1000}$.  However, we still have to make sure a sufficiently small $r$ exists.  

 In light of the equation $m - n^3 = r(3n^2 + 3nr + r^2)$, we need to choose $m - n^3$ as small as possible to ensure a small enough $r$.  The smallest possible value for $m - n^3$ is 1, when $m = 19^3 + 1$.  Then for this value of $m$, $r = \frac{1}{3n^2 + 3nr + r^2} < \frac{1}{1000}$, and we're set.  The answer is  $\boxed{019}$.

 



\textbf{Solution 2}

To minimize $m$, we should minimize $n$. We have that $(n + \frac{1}{1000})^3 = n^3 + \frac{3}{10^3} n^2 + \frac{3}{10^6} n + \frac{1}{10^9}$. For a given value of $n$, if $(n + \frac{1}{1000})^3 - n^3 > 1$, there exists an integer between $(n + \frac{1}{1000})^3$ and $n^3$, and the cube root of this integer would be between $n$ and $n + \frac{1}{1000}$ as desired. We seek the smallest $n$ such that $(n + \frac{1}{1000})^3 - n^3 > 1$. 

 
 \[(n + \frac{1}{1000})^3 - n^3 > 1\]
 \[\frac{3}{10^3} n^2 + \frac{3}{10^6} n + \frac{1}{10^9} > 1\]
 \[3n^2 + \frac{3}{10^3} n + \frac{1}{10^6} > 10^3\]

 Trying values of $n$, we see that the smallest value of $n$ that works is $\boxed{019}$.

 Why is it $(n + \frac{1}{1000})^3 - n^3 > 1$ and not greater than or equal to? - awesomediabrine

 Because if its equal to, then there is no integer in between the two values. - resources

 



\textbf{Solution 3 (Similar to Solution 2)}

Since $r$ is less than $1/1000$, we have $\sqrt[3]{m} < n + \frac{1}{1000}$. Notice that since we want $m$ minimized, $n$ should also be minimized. Also, $n^3$ should be as close as possible, but not exceeding $m$. This means $m$ should be set to $n^3+1$. Substituting and simplifying, we get 
 \[\sqrt[3]{n^3+1} < n + \frac{1}{1000}\] 
 \[n^3+1 < n^3+\frac{3}{1000}n^2+\frac{3}{1000^2}n+\frac{1}{1000^3}\]The last two terms in the right side can be ignored in the calculation because they are too small. This results in $1 < \frac{3}{1000}n^2 \Rightarrow n^2 > \frac{1000}{3}$. The minimum positive integer $n$ that satisfies this is $\boxed{019}$. ~ Hb10

 



\textbf{Solution 4 (Calculus)}

Note that the cube root is increasing for positive reals while its derivative is decreasing, so linear approximation gives
 \[\sqrt[3]{n^3+1} - n < \left.\frac{d\sqrt[3]{x}}{dx}\right|_{x=n^3} = \frac{1}{3n^2}\]and
 \[\sqrt[3]{n^3+1} - n > \left.\frac{d\sqrt[3]{x}}{dx}\right|_{x=n^3+1} = \frac{1}{3\sqrt[3]{(n^3+1)^2}}\]From this, it is clear that $n = \boxed{019}$ is the smallest $n$ for which LHS will be less than $\frac{1}{1000}$. ~ Hyprox1413

 


	\item (\textit{AIME I 1987}) A given sequence $r_1, r_2, \dots, r_n$ of distinct real numbers can be put in ascending order by means of one or more "bubble passes".  A bubble pass through a given sequence consists of comparing the second term with the first term, and exchanging them if and only if the second term is smaller, then comparing the third term with the second term and exchanging them if and only if the third term is smaller, and so on in order, through comparing the last term, $r_n$, with its current predecessor and exchanging them if and only if the last term is smaller. 


 The example below shows how the sequence 1, 9, 8, 7 is transformed into the sequence 1, 8, 7, 9 by one bubble pass.  The numbers compared at each step are underlined.


 $\underline{1 \quad 9} \quad 8 \quad 7$
$1 \quad {}\underline{9 \quad 8} \quad 7$
$1 \quad 8 \quad \underline{9 \quad 7}$
$1 \quad 8 \quad 7 \quad 9$

 Suppose that $n = 40$, and that the terms of the initial sequence $r_1, r_2, \dots, r_{40}$ are distinct from one another and are in random order.  Let $p/q$, in lowest terms, be the probability that the number that begins as $r_{20}$ will end up, after one bubble pass, in the $30^{\mbox{th}}$ place.  Find $p + q$.


 
 \textbf{Answer: }931



\textbf{Solution 1}

If any of $r_1, \ldots, r_{19}$ is larger than $r_{20}$, one of these numbers will be compared with $r_{20}$ on the 19th step of the first bubble pass and $r_{20}$ will be moved back to the 19th position.  Thus, $r_{20}$ must be the largest of the first 20 terms.  In addition, $r_{20}$ must be larger than $r_{21}, r_{22}, \ldots, r_{30}$ but smaller than $r_{31}$ in order that it move right to the 30th position but then not continue moving right to the 31st.

 Thus, our problem can be restated: What is the probability that in a sequence of 31 distinct real numbers, the largest is in position 31 and the second-largest is in position 20 (the other 29 numbers are irrelevant)?

 This is much easier to solve: there are $31!$ ways to order the first thirty-one numbers and $29!$ ways to arrange them so that the largest number is in the 31st position and the second-largest is in the 20th. This gives us a desired probability of $\frac{29!}{31!} = \frac{1}{31\cdot 30} = \frac{1}{930}$, so the answer is $\boxed{931}$.

 \textbf{Remark} Note that you can solve the restated problem differently: 

 We see that the numerator is $1$ because there is only $1$ way to make the $31^{\text{st}}$ term the largest and the $20^{\text{th}}$ the second-largest. To calculate the denominator, we note that there are $31\times30$ ways to pick the largest term and the second-largest term. Therefore, our answer is $\boxed{931}$.~Yiyj1

 



\textbf{Solution 2 (Summations)}

Assume we have a set of $40$ numbers, all distinct, that will be randomly shuffled. We wish to find the number of ways that the following conditions are satisfied:

 $1.$ The greatest number out of the first $30$ terms is located at the $20$th place, and

 $2.$ Term $31$ is greater than term $20$.

 $3.$ The final $9$ terms can be in any order.

 Let the resulting term $20$ be notated as $X$. We then set up a dummy variable $k$ such that $k$ describes the number of terms (out of $40$) that are less than $X$. For instance, if there were $29$ terms less than $X$, the dummy variable $k$ would be $29$. Now we can set up our summation.

 $1.$ To make sure the first condition is met, we perform the following: We know that there are $k$ terms less than $X$, so thus there are $\frac{k!}{(k-29)!}$ ways to place $29$ of the $k$ terms in a specific order.

 $2.$ For the second condition, there are $39-k$ terms greater than $X$, so there are $39-k$ ways to place the number greater than $X$ in the $31$st placement.

 $3.$ Finally, since the final $9$ terms can be in any order, there are $9!$ ways to place these numbers.

 (The three steps above all correspond to the three conditions given at the beginning of the solution.)

 However, we need to set lower and upper bounds for $k$ in order to calculate our summation. Looking at the range of $k$, we see that there must be at least $29$ numbers less than $X$ in order for condition $1$ to be satisfied. There must be at least $1$ number greater than $X$ in order for condition $2$ to be satisfied as well, leaving a maximum of $38$ numbers less than $X$. Thus, $k$ ranges from $29$ to $38$.

 Finally, we divide the entire value by $40!$ as there are $40!$ ways to arrange the $40$ terms. This step could also be performed at the very end. At this point, we are ready to calculate the summation.

 
$\frac{1}{40!}\sum_{k=29}^{38}\left(\frac{k!}{(k-29)!}(39-k)\cdot9!\right)$
$=\frac{9!}{40!}\sum_{k=29}^{38}\left(\frac{k!}{(k-29)!}(39-k)\right)$
$=\frac{9!\cdot39}{40!}\sum_{k=29}^{38}\left(\frac{k!}{(k-29)!}\right)-\frac{9!}{40!}\sum_{k=29}^{38}\left(\frac{k!\cdot k}{(k-29)!}\right)$

 Now we deal with each sum separately. Looking at the first sum, we see that the factorials inside the summation bear a close resemblance to ${k\choose29}$. Indeed, we are able to use this to our advantage:

 $\frac{9!\cdot39}{40!}\sum_{k=29}^{38}\left(\frac{k!}{(k-29)!}\right)$
$=\frac{9!\cdot39\cdot29!}{40!}\sum_{k=29}^{38}\left(\frac{k!}{(k-29)!\cdot29!}\right)$
$=\frac{9!\cdot39\cdot29!}{40!}\sum_{k=29}^{38}{k\choose29}$
$=\frac{9!\cdot39\cdot29!}{40!}{39\choose30}$

 (The last equality is thanks to the \href{https://artofproblemsolving.com/wiki/index.php?title=Hockey-Stick\_Identity}{Hockey-Stick Identity}.)

 The second sum is similar to the first, but we must first replace $k$ with $(k+1)-1$:

 $\frac{9!}{40!}\sum_{k=29}^{38}\left(\frac{k!\cdot k}{(k-29)!}\right)$
$=\frac{9!}{40!}\sum_{k=29}^{38}\left(\frac{k!\cdot(k+1-1)}{(k-29)!}\right)$
$=\frac{9!}{40!}\sum_{k=29}^{38}\left(\frac{k!\cdot(k+1)}{(k-29)!}\right)-\frac{9!}{40!}\sum_{k=29}^{38}\left(\frac{k!}{(k-29)!}\right)$
$=\frac{9!}{40!}\sum_{k=29}^{38}\left(\frac{(k+1)!}{(k-29)!}\right)-\frac{9!}{40!}\sum_{k=29}^{38}\left(\frac{k!}{(k-29)!}\right)$
$=\frac{9!\cdot30!}{40!}\sum_{k=29}^{38}{k+1\choose30}-\frac{9!\cdot29!}{40!}\sum_{k=29}^{38}{k\choose29}$
$=\frac{9!\cdot30!}{40!}{40\choose31}-\frac{9!\cdot29!}{40!}{39\choose30}$

 We are finally able to calculate the end result:

 $\frac{9!\cdot39}{40!}\sum_{k=29}^{38}\left(\frac{k!}{(k-29)!}\right)-\frac{9!}{40!}\sum_{k=29}^{38}\left(\frac{k!\cdot k}{(k-29)!}\right)$
$=\frac{9!\cdot39\cdot29!}{40!}{39\choose30}-\frac{9!\cdot30!}{40!}{40\choose31}+\frac{9!\cdot29!}{40!}{39\choose30}$
$=\frac{9!\cdot40\cdot29!}{40!}{39\choose30}-\frac{9!\cdot30!}{40!}{40\choose31}$
$=\frac{9!\cdot29!\cdot39!}{39!\cdot9!\cdot30!}-\frac{9!\cdot30!\cdot40!}{40!\cdot9!\cdot31!}$
$=\frac{29!}{30!}-\frac{30!}{31!}$
$=\frac{1}{30}-\frac{1}{31}$
$=\frac{31}{930}-\frac{30}{930}$
$=\frac{1}{930}$

 Thus our answer is $1+930=\boxed{931}$.

 


	\item (\textit{AIME I 1987}) Compute

 \[\frac{(10^4+324)(22^4+324)(34^4+324)(46^4+324)(58^4+324)}{(4^4+324)(16^4+324)(28^4+324)(40^4+324)(52^4+324)}.\]

 \textbf{Answer: }373



\textbf{Solution 1 (Sophie Germain Identity)}

The \href{https://artofproblemsolving.com/wiki/index.php?title=Sophie\_Germain\_Identity}{Sophie Germain Identity} states that $a^4 + 4b^4$ can be factored as $\left(a^2 + 2b^2 - 2ab\right)\left(a^2 + 2b^2 + 2ab\right).$ Each of the terms is in the form of $x^4 + 324.$ Using Sophie Germain, we get that
 \begin{align*} x^4 + 324 &= x^4 + 4\cdot 3^4 \\ &= \left(x^2 + 2 \cdot 3^2 - 2\cdot 3\cdot x\right)\left(x^2 + 2 \cdot 3^2 + 2\cdot 3\cdot x\right) \\ &= (x(x-6) + 18)(x(x+6)+18), \end{align*}
so the original expression becomes
 \[\frac{[(10(10-6)+18)(10(10+6)+18)][(22(22-6)+18)(22(22+6)+18)]\cdots[(58(58-6)+18)(58(58+6)+18)]}{[(4(4-6)+18)(4(4+6)+18)][(16(16-6)+18)(16(16+6)+18)]\cdots[(52(52-6)+18)(52(52+6)+18)]},\] which simplifies to 
 \[\frac{(10(4)+18)(10(16)+18)(22(16)+18)(22(28)+18)\cdots(58(52)+18)(58(64)+18)}{(4(-2)+18)(4(10)+18)(16(10)+18)(16(22)+18)\cdots(52(46)+18)(52(58)+18)}.\]Almost all of the terms cancel out! We are left with $\frac{58(64)+18}{4(-2)+18} = \frac{3730}{10} = \boxed{373}.$




\textbf{Solution 2 (Completing the Square and Difference of Squares)}

In both the numerator and the denominator, each factor is of the form $N^4+324=N^4+18^2$ for some positive integer $N.$ 

 We factor $N^4+18^2$ by completing the square, then applying the difference of squares:
 \begin{align*} N^4+18^2&=\left(N^4+36N^2+18^2\right)-36N^2 \\ &=\left(N^2+18\right)^2-(6N)^2 \\ &=\left(N^2-6N+18\right)\left(N^2+6N+18\right) \\ &=\left((N-3)^2+9\right)\left((N+3)^2+9\right). \end{align*}
The original expression now becomes 
 \[\frac{\left[(7^2+9)(13^2+9)\right]\left[(19^2+9)(25^2+9)\right]\left[(31^2+9)(37^2+9)\right]\left[(43^2+9)(49^2+9)\right]\left[(55^2+9)(61^2+9)\right]}{\left[(1^2+9)(7^2+9)\right]\left[(13^2+9)(19^2+9)\right]\left[(25^2+9)(31^2+9)\right]\left[(37^2+9)(43^2+9)\right]\left[(49^2+9)(55^2+9)\right]}=\frac{61^2+9}{1^2+9}=\boxed{373}.\]~MRENTHUSIASM

 



\textbf{Solution 3 (Complex Numbers)}

In both the numerator and the denominator, each factor is of the form $N^4+324=N^4+18^2$ for some positive integer $N.$ 

 We factor $N^4+18^2$ by solving the equation $N^4+18^2=0,$ or $N^4=-18^2.$

 Two solutions follow from here:

 



\textbf{Solution 3.1 (Polar Form)}

We rewrite $N$ to the polar form 
 \[N=r(\cos\theta+i\sin\theta)=r\operatorname{cis}\theta,\] where $r$ is the magnitude of $N$ such that $r\geq0,$ and $\theta$ is the argument of $N$ such that $0\leq\theta<2\pi.$

 By \textbf{De Moivre's Theorem}, we have 
 \[N^4=r^4\operatorname{cis}(4\theta)=18^2(-1),\] from which

 
            \begin{enumerate} 
        \item $r^4=18^2,$ so $r=3\sqrt2.$

        \item $\begin{cases} \begin{aligned} \cos(4\theta) &= -1 \\ \sin(4\theta) &= 0 \end{aligned}, \end{cases}$ so $\theta=\frac{\pi}{4},\frac{3\pi}{4},\frac{5\pi}{4},\frac{7\pi}{4}.$

            \end{enumerate}
            By the \textbf{Factor Theorem}, we get
 \begin{align*} N^4+18^2&=\biggl(N-3\sqrt2\operatorname{cis}\frac{\pi}{4}\biggr)\biggl(N-3\sqrt2\operatorname{cis}\frac{3\pi}{4}\biggr)\biggl(N-3\sqrt2\operatorname{cis}\frac{5\pi}{4}\biggr)\biggl(N-3\sqrt2\operatorname{cis}\frac{7\pi}{4}\biggr) \\ &=\biggl[\biggl(N-3\sqrt2\operatorname{cis}\frac{\pi}{4}\biggr)\biggl(N-3\sqrt2\operatorname{cis}\frac{7\pi}{4}\biggr)\biggr]\biggl[\biggl(N-3\sqrt2\operatorname{cis}\frac{3\pi}{4}\biggr)\biggl(N-3\sqrt2\operatorname{cis}\frac{5\pi}{4}\biggr)\biggr] \\ &=\left[(N-(3+3i))(N-(3-3i))\right]\left[(N-(-3+3i))(N-(-3-3i))\right] \\ &=\left[((N-3)-3i)((N-3)+3i)\right]\left[((N+3)-3i)((N+3)+3i)\right] \\ &=\left[(N-3)^2+9\right]\left[(N+3)^2+9\right]. \end{align*}
We continue with the last paragraph of Solution 2 to get the answer $\boxed{373}.$




\textbf{Solution 3.2 (Rectangular Form)}

We rewrite $N$ to the rectangular form 
 \[N=a+bi\] for some real numbers $a$ and $b.$

 Note that $N^2=\pm18i,$ so there are two cases:

 
            \begin{enumerate} 
        \item $N^2=18i$We have
 \begin{align*} (a+bi)^2&=18i \\ a^2-b^2+2abi&=18i. \end{align*}
We need $\begin{cases} \begin{aligned} a^2-b^2 &= 0 \\ 2ab &= 18 \end{aligned}, \end{cases}$ from which $(a,b)=(3,3),(-3,-3),$ or $N=3+3i,-3-3i.$


        \item $N^2=-18i$We have
 \begin{align*} (a+bi)^2&=-18i \\ a^2-b^2+2abi&=-18i. \end{align*}
We need $\begin{cases} \begin{aligned} a^2-b^2 &= 0 \\ 2ab &= -18 \end{aligned}, \end{cases}$ from which $(a,b)=(3,-3),(-3,3),$ or $N=3-3i,-3+3i.$


            \end{enumerate}
            By the \textbf{Factor Theorem}, we get
 \begin{align*} N^4+18^2&=(N-(3+3i))(N-(-3-3i))(N-(3-3i))(N-(-3+3i)) \\ &=\left[(N-(3+3i))(N-(3-3i))\right]\left[(N-(-3+3i))(N-(-3-3i))\right] \\ &=\left[((N-3)-3i)((N-3)+3i)\right]\left[((N+3)-3i)((N+3)+3i)\right] \\ &=\left[(N-3)^2+9\right]\left[(N+3)^2+9\right]. \end{align*}
We continue with the last paragraph of Solution 2 to get the answer $\boxed{373}.$




\textbf{Solution 4}

We use Sophie Germain's Identity to rewrite the first couple of multiplicands in the numerator and denominator. By Sophie Germain's: 
 \[(10^4 + 324) = (10^4 + \cdot 4 \cdot 3^4) = (10^2 + 2 \cdot 10 \cdot 3 + 2 \cdot 3^2)(10^2 - 2 \cdot 10 \cdot 3 + 2 \cdot 3^2) = (178)(58)\] 
 \[(22^4 + 324) = (22^2 + 2 \cdot 3 \cdot 22 + 2 \cdot 3^2)(22^2 - 2 \cdot 3 \cdot 22 + 2 \cdot 3^2) = (634)(370)\] 
 \[(4^4 + 324) = (4^2 + 2 \cdot 4 \cdot 3 + 2 \cdot 3^2)(4^2 - 2 \cdot 4 \cdot 3 + 2 \cdot 3^2) = (58)(10)\] 
 \[(16^4 + 324) = (16^2 + 2 \cdot 16 \cdot 3 + 2 \cdot 3^2)(16^2 - 2 \cdot 16 \cdot 3 + 2 \cdot 3^2) = (370)(178)\] If we only had these terms, then the fraction would rewrite to $\frac{(58)(178)(370)(634)}{(10)(58)(178)(370)}$. However, we notice most of the terms cancel, leaving us only with the largest term in the numerator and the smallest term in the denominator ($\frac{634}{10}$). We hypothesize that this will happen with the fraction as a whole. Then we will only be left with the largest term in the numerator, which is $(58^2 + 2 \cdot 58 \cdot 3 + 2 \cdot 3^2) = 3730$. The fraction simplifies to $\frac{3730}{10} = \boxed{373}$.

 


	\item (\textit{AIME I 1987}) Squares $S_1$ and $S_2$ are inscribed in right triangle $ABC$, as shown in the figures below. Find $AC + CB$ if area $(S_1) = 441$ and area $(S_2) = 440$.


 
                    \begin{figure}[H]
                        \centering
                        \includegraphics[width=0.5\textwidth]{Problem/AIME_I_1987_15_Problem_Figure_1.png}
                    \end{figure}
                    

 \textbf{Answer: }462



\textbf{Solution}


                    \begin{figure}[H]
                        \centering
                        \includegraphics[width=0.5\textwidth]{Solution/AIME_I_1987_15_Solution_Figure_1.png}
                    \end{figure}
                    

 Because all the \href{https://artofproblemsolving.com/wiki/index.php?title=Triangle}{triangles} in the figure are \href{https://artofproblemsolving.com/wiki/index.php?title=Similar}{similar} to triangle $ABC$, it's a good idea to use \href{https://artofproblemsolving.com/wiki/index.php?title=Area\_ratios}{area ratios}. In the diagram above, $\frac {T_1}{T_3} = \frac {T_2}{T_4} = \frac {441}{440}.$ Hence, $T_3 = \frac {440}{441}T_1$ and $T_4 = \frac {440}{441}T_2$. Additionally, the area of triangle $ABC$ is equal to both $T_1 + T_2 + 441$ and $T_3 + T_4 + T_5 + 440.$ 

 Setting the equations equal and solving for $T_5$, $T_5 = 1 + T_1 - T_3 + T_2 - T_4 = 1 + \frac {T_1}{441} + \frac {T_2}{441}$. Therefore, $441T_5 = 441 + T_1 + T_2$. However, $441 + T_1 + T_2$ is equal to the area of triangle $ABC$! This means that the ratio between the areas $T_5$ and $ABC$ is $441$, and the ratio between the sides is $\sqrt {441} = 21$. As a result, $AB = 21\sqrt {440} = \sqrt {AC^2 + BC^2}$. We now need $(AC)(BC)$ to find the value of $AC + BC$, because $AC^2 + BC^2 + 2(AC)(BC) = (AC + BC)^2$.

 Let $h$ denote the height to the \href{https://artofproblemsolving.com/wiki/index.php?title=Hypotenuse}{hypotenuse} of triangle $ABC$. Notice that $h - \frac {1}{21}h = \sqrt {440}$. (The height of $ABC$ decreased by the corresponding height of $T_5$) Thus, $(AB)(h) = (AC)(BC) = 22\cdot 21^2$. Because $AC^2 + BC^2 +  2(AC)(BC) = (AC + BC)^2 = 21^2\cdot22^2$, $AC + BC = (21)(22) = \boxed{462}$.

 \textbf{Easy Trig Solution} Let $\tan\angle ABC = x$. Now using the 1st square, $AC=21(1+x)$ and $CB=21(1+x^{-1})$. Using the second square, $AB=\sqrt{440}(1+x+x^{-1})$. We have $AC^2+CB^2=AB^2$, or 
 \[441(x^2+x^{-2}+2x+2x^{-1}+2)=440(x^2+x^{-2}+2x+2x^{-1}+3).\] Rearranging and letting $u=x+x^{-1} \Rightarrow u^2 - 2 = x^2 + x^{-2}$ gives us $u^2+2u-440=0.$ We take the positive root, so $u=20$, which means $AC+CB=21(2+x+x^{-1})=21(2+u)=\boxed{462}$.

 \textbf{Messy Trig Solution} Let $\theta$ be the smaller angle in the triangle. Then the sum of shorter and longer leg is $\sqrt{441}(2+\tan{\theta}+\cot{\theta})$. We observe that the short leg has length $\sqrt{441}(1+\tan{\theta}) = \sqrt{440}(\sec{\theta}+\sin{\theta})$. Grouping and squaring, we get $\sqrt{\frac{440}{441}} = \frac{\sin{\theta}+\cos{\theta}}{1+\sin{\theta}\cos{\theta}}$. Squaring and using the double angle identity for sine, we get, $110(\sin{2\theta})^2 + \sin{2\theta} - 1 = 0$. Solving, we get $\sin{2\theta} = \frac{1}{10}$. Now to find $\tan{\theta}$, we find $\cos{2\theta}$ using the PythagoreanIdentity, and then use the tangent double angle identity. Thus, $\tan{\theta} = 10-3\sqrt{11}$. Substituting into the original sum,we get $\boxed{462}$.

 



\textbf{Solution 4 (Algebra)}


                    \begin{figure}[H]
                        \centering
                        \includegraphics[width=0.5\textwidth]{Solution/AIME_I_1987_15_Solution_Figure_2.png}
                    \end{figure}
                    

                    \begin{figure}[H]
                        \centering
                        \includegraphics[width=0.5\textwidth]{Solution/AIME_I_1987_15_Solution_Figure_3.png}
                    \end{figure}
                    

 Label points as above. Let $x=AC$, $y=BC$, $s_1 = 21$ be the side length of $S_1$, and $s_2 = \sqrt{440}$ be the side length of $S_2$.

 Since $\triangle ABC\sim\triangle AED$, we have $\frac{x}{y} = \frac{x-s_1}{s_1}$
$\implies xs_1=xy-ys_1$
$\implies xy=s_1(x+y)$
$\implies xy=21(x+y) \qquad \qquad (*)$.

 Since $\triangle ABC\sim\triangle AWX\sim\triangle ZBY$, we have $s_2 + \frac{s_2x}{y} + \frac{s_2y}{x}=\sqrt{x^2+y^2}$
$\implies s_2(x^2+xy+y^2)=xy\sqrt{x^2+y^2}$
$\implies s_2^2(x^2+xy+y^2)^2 = x^2y^2(x^2+y^2)$
$\implies 440(x^2+xy+y^2)^2 = x^2y^2(x^2+y^2)$

 Let $t=x+y$. Repeatedly applying $(*)$, we get
 \[440(t^2-21t)^2 = 441t^2(t^2 - 42t)\]
 \[440(t-21)^2 = 441(t^2-42t)\]
 \[440t^2 - 42\cdot 440t + 440\cdot 441 = 441t^2 - 441\cdot 42t\]
 \[t^2-42t-440\cdot 441=0\]
 \[(t-21)^2 = 441^2\]
 \[t-21=441\]
 \[t=\boxed{462}\]



\end{enumerate}